\documentclass[11pt,a4paper]{article}
\usepackage{amsmath,amssymb,amsthm}
\usepackage{hyperref}
\hypersetup{
  pdftitle={Temporal Residual Map from the Completed Principal Symbol: Construction, Uniqueness, and the Lorentz Identification},
  pdfauthor={J\'er\^ome Beau}
}
\usepackage{doi}
\usepackage{geometry}
\geometry{margin=2.5cm}

\newtheorem{theorem}{Theorem}[section]
\newtheorem{lemma}[theorem]{Lemma}
\newtheorem{corollary}[theorem]{Corollary}
\newtheorem{proposition}[theorem]{Proposition}
\newtheorem{definition}[theorem]{Definition}
\newtheorem{remark}[theorem]{Remark}

\newcommand{\Sym}{\mathrm{Sym}}

\title{Temporal Residual Map from the Completed Principal Symbol:\\
Construction, Uniqueness, and the Lorentz Identification}
\author{J\'er\^ome Beau\\
\small Independent Researcher\\
\small \texttt{jerome.beau@cosmochrony.org}}
\date{2026}

\begin{document}

  \maketitle

  \begin{abstract}
    The LorCap note identifies the construction of a temporal residual
    map~$\mathcal{P}_\tau^{\mathrm{Q5b}}$ as the remaining step needed to close the Lorentz
    identification from O7 observables.
    The present paper carries out this construction from three inputs:
    the temporal splitting of a supplied rank-four co-metric of the conditional form of Q11,
    $\mathrm{diag}(-A_\tau, 2\lambda, 2\lambda, 2\lambda)$ with $A_\tau$ and $\lambda$ independent and fixed by no
    source,
    the Born--Infeld unit capacity-flow sphere (the quadratic budget hypothesis expressing
    the Born--Infeld capacity bound), and
    the exact Weil linear filling law $\Sigma_n^{\mathrm{tot}} = 1 - B_n$
    established in LorCap, with $B_n$ identified with the spatial capacity flow by a hypothesis of the paper.
    The splitting uses only the absence of mixed temporal--spatial terms and a non-zero temporal coefficient; neither a coefficient value nor the
    signature enters the construction, and the normalisation $g^{\mu\nu} = 2\eta^{\mu\nu}$ is a choice of units.
    The resulting scalar temporal residual map is
    \[
      \mathcal{P}_\tau^{\mathrm{Q5b}} \colon \Sigma_n^{\mathrm{tot}}
      \;\longmapsto\; \sqrt{\Sigma_n^{\mathrm{tot}}\bigl(2 - \Sigma_n^{\mathrm{tot}}\bigr)}
      \;=\; \sqrt{1 - B_n^2}.
    \]
    It is not an idempotent projection on scalar observables; it is the scalar shadow
    of the linear temporal projector of the co-metric composed with the Born--Infeld lift.

    Three statements are distinct.
    First, scalar uniqueness: the quadratic law and the Weil filling law alone force the values of any
    non-negative map at the values of $\Sigma_n^{\mathrm{tot}}$ actually attained, a finite set of points in
    $[1/2 - 1/(2q^2),\, 1 - 1/q^2]$ for depths below half the torus; off that set the map is not constrained.
    Second, the splitting identifies that value as the temporal component of the capacity flow; it does not
    select the formula, and the uniqueness does not use it.
    Third, the Lorentzian kinematic reading uses the signature of the co-metric and a further hypothesis,
    [H-cone]: the light-cone speed of the co-metric equals the Born--Infeld bound per ordering step.
    This is the coupling datum between temporal and spatial normalisations that Q11 leaves open; it is not derived
    here.
    The three candidate maps tested numerically in LorCap all fail the quadratic law; the paper states how each
    differs from the construction.
    LorCap Lemma~1 follows, and with it LorCap Theorem~1, as the exact identity
    $\Phi(\eta_n^\tau) = (1+(\eta_n^\tau)^2)^{-1/2} = (1-B_n^2)^{1/2}$.
    \textbf{Interpretive reading.}
    The Lorentz factor is thereby read as the reciprocal temporal residual of a bounded
    capacity-flow split, conditional on the supplied co-metric, the identification of $B_n$, the Born--Infeld lift
    and, for its kinematic reading, [H-cone].
  \end{abstract}

  \noindent\textbf{Keywords.}
  temporal residual map, principal symbol, Born--Infeld admissibility, Lorentz factor,
  capacity-flow bundle, uniqueness, Heisenberg group, sub-Riemannian geometry,
  effective co-metric, Weil factorization, Cosmochrony.

  \newpage
  \tableofcontents

  \newpage

  \section{Introduction}

    \subsection{Context and the open problem}

      The Cosmochrony spectral geometry programme seeks to derive effective spacetime structure from
      admissibility constraints on non-injective projections of spectral data on the Heisenberg
      group $\mathrm{Heis}_3(\mathbb{Z}/q\mathbb{Z})$~\cite{Beau2026m}.
      Under the spatial-limit hypothesis [H-L] and the lifting hypothesis [H-lift], the paper
      Q5b~\cite{Beau2026q5b} obtains the effective operator
      $L_{\mathrm{eff}} = -A_\tau\partial_\tau^2 + A_H\Delta_H$ on
      $\mathbb{R}_\tau \times \mathrm{Heis}_3(\mathbb{R})$, whose principal symbol gives a co-metric of rank three;
      its rank-four extension (open problem Q5b-O2) and the coefficient values (Q5b-O3) are open.
      No source fixes the coefficients.
      Q8~\cite{BeauQ8} shows that a supplied, fully invariant spatial form has $A_H = A_Z = 2\lambda$ with
      $\lambda > 0$ free, Q10~\cite{Beau2026q10} derives no coefficient value, and Q11~\cite{Beau2026q11} shows that
      spatial $\mathrm{SU}(2)$ invariance leaves $A_\tau$ independent of $\lambda$.
      Under the hypotheses of Q11 Corollary~6.1 the co-metric is therefore
      $\mathrm{diag}(-A_\tau, 2\lambda, 2\lambda, 2\lambda)$ with two independent scales.

      The note~\cite{Beau2026n} then studies the O7 saturation observable $\eta_n^{O7} =
  B_n/\Sigma_n^{\mathrm{tot}}$ and asks whether it recovers the relativistic Lorentz factor
      $\gamma$.
      Two structural facts are established there:
      \begin{enumerate}
        \item The Weil factorization gives $\Sigma_n^{\mathrm{tot}} = 1 - B_n$ exactly in the
        Schr\"odinger representation (Corollary~1 of~\cite{Beau2026n}).
        \item A temporal residual projection $\mathcal{P}_\tau^{\mathrm{Q5b}}$ is needed;
        no natural BFS observable directly satisfies the quadratic law
        $B_n^2 + (\Sigma_n^\tau)^2 = 1$.
      \end{enumerate}
      The note defines $\mathcal{P}_\tau^{\mathrm{Q5b}}$ implicitly via a completed
      co-metric, but defers its explicit construction as ``the remaining open
      task'' (Section~6.3 of~\cite{Beau2026n}).

    \subsection{What the present paper does}

      We carry out this construction in full (Theorem~\ref{thm:main}) and separate three statements.
      The scalar uniqueness follows from the quadratic BI law and the Weil filling law alone: they force the
      values of any non-negative map at the attained values of $\Sigma_n^{\mathrm{tot}}$, and only there.
      The temporal splitting of the co-metric identifies that value as the temporal component of the capacity
      flow; it does not select the formula.
      The Lorentzian kinematic reading uses in addition the signature of the co-metric and the hypothesis
      [H-cone] (Corollary~\ref{cor:lorentz}).

      The three candidate maps $\Sigma_n^A$, $\Sigma_n^B$, $\Sigma_n^C$ tested numerically in~\cite{Beau2026n}
      all fail the quadratic law; Proposition~\ref{prop:failure} states how each differs from the construction.

    \subsection{Precise scope of each ingredient}\label{sec:scope}

      The construction rests on three distinct inputs, which play different roles.
      \begin{enumerate}
        \item[(L)] \textbf{Temporal splitting} (from a supplied completed principal symbol):
        a co-metric of the form~\eqref{eq:metric} determines the timelike direction
        $\partial_\tau$ in $T(\mathbb{R}_\tau \times \mathrm{Heis}_3(\mathbb{R}))$, orthogonal to the spatial
        directions because its temporal block is decoupled, and a corresponding linear projector~$P_\tau$ on
        vectors.
        The projector uses the absence of mixed terms, not the values of $A_\tau$ and $\lambda$; the signature
        makes $\partial_\tau$ timelike but enters only the kinematic reading (Corollary~\ref{cor:lorentz}).
        (L) identifies which component of $\vec{F}_n$ is temporal; it does not constrain the scalar formula.
        The principal symbol alone does not determine a scalar map on $\Sigma_n^{\mathrm{tot}}$.
        \item[(BI)] \textbf{Born--Infeld capacity sphere} (from the quadratic admissibility
        budget expressing the capacity bound of~\cite{Beau2026c}): the admissible capacity flow vector
        $\vec{F}_n$ lies on the unit sphere $(F_n^\tau)^2 + |\vec{F}_n^{\,\mathrm{spat}}|^2 = 1$ (after normalisation).
        This is the quadratic constraint that forces $F_n^\tau = \sqrt{1 - B_n^2}$ once $B_n$
        is identified with the spatial component (identification (S), Section~\ref{sec:spatialid}).
        \item[(W)] \textbf{Weil filling law} (from the Schr\"odinger factorization~\cite{Beau2026n}):
        $\Sigma_n^{\mathrm{tot}} = 1 - B_n$ exactly, expressing $B_n$ through the observable as
        $B_n = 1 - \Sigma_n^{\mathrm{tot}}$.
      \end{enumerate}
      The quadratic law (BI) and the identity (W) determine the values of a non-negative map at the attained
      values of $\Sigma_n^{\mathrm{tot}}$ (Theorem~\ref{thm:main}(ii)); (L) identifies that value with the temporal
      component of the capacity flow (Theorem~\ref{thm:main}(iii)).

    \subsection{Organisation}

      Section~\ref{sec:setup} fixes notation and recalls the relevant results.
      Section~\ref{sec:bundle} defines the capacity-flow bundle and the unit BI sphere.
      Section~\ref{sec:construction} constructs $\mathcal{P}_\tau^{\mathrm{Q5b}}$ and proves
      Theorem~\ref{thm:main}.
      Section~\ref{sec:necessity} states which inputs the uniqueness needs and compares the
      three candidates of~\cite{Beau2026n} with the construction.
      Section~\ref{sec:lorentz} derives the Lorentz identification.
      Section~\ref{sec:status} gives the logical status table.

  \section{Setup and Notation}\label{sec:setup}

    \subsection{Objects in force}\label{sec:objects}

      Throughout, $q$ denotes a prime, $G_q = \mathrm{Heis}_3(\mathbb{Z}/q\mathbb{Z})$,
      and $d_w$ denotes the BFS word metric on $G_q$.
      The O7 cumulative occupation observable at depth $n < q/2$ (proved exact
      in~\cite{Beau2026n}) is
      \begin{equation}\label{eq:Bn}
        B_n \;:=\; \frac{2n^2+2n+1}{q^2}, \qquad n < q/2.
      \end{equation}
      This is the fraction of $(a,b)$ pairs in $({\mathbb{Z}/q\mathbb{Z}})^2$ satisfying
      $\min(a,q{-}a)+\min(b,q{-}b)\leq n$, i.e.\ the normalised spatial footprint of the BFS
      ball at depth~$n$.
      Let $R := \lfloor (q-1)/2 \rfloor$, the largest depth below $q/2$.
      Then $B_0 = 1/q^2$, $B_n$ is strictly increasing for $0 \le n \le R$, and for odd $q$ its largest value is
      $B_R = (q^2+1)/(2q^2) = 1/2 + 1/(2q^2)$, slightly above $1/2$.
      The formula~\eqref{eq:Bn} does not extend beyond $n = q/2$; the full saturation profile
      on the finite torus involves wrapping corrections.
      The O7 total residual frontier novelty is the exact complement (Corollary~1
      of~\cite{Beau2026n}):
      \begin{equation}\label{eq:Sigma}
        \Sigma_n^{\mathrm{tot}} \;=\; 1 - B_n.
      \end{equation}
      In the window $0 \le n \le R$ the total residual therefore takes exactly the $R+1$ distinct values
      \[
        S_q := \{\Sigma_n^{\mathrm{tot}} : 0 \le n \le R\}
        \;\subset\; \bigl[\tfrac12 - \tfrac{1}{2q^2},\; 1 - \tfrac{1}{q^2}\bigr]
      \]
      (for odd $q$), the \emph{attained values}; for $q = 11$ there are six.

    \subsection{Completed co-metric and its inverse}\label{sec:cometric}

      Under [H-L] and [H-lift], the effective operator on $\mathbb{R}_\tau \times \mathrm{Heis}_3(\mathbb{R})$ is
      (Q5b Theorem~5.2~\cite{Beau2026q5b}):
      \begin{equation}\label{eq:Leff}
        L_{\mathrm{eff}} = -A_\tau\partial_\tau^2 + A_H\Delta_H,
      \end{equation}
      where $\Delta_H = \tilde{X}^2 + \tilde{Y}^2$ is the Kohn--Laplacian on $\mathrm{Heis}_3(\mathbb{R})$.
      The principal symbol of $L_{\mathrm{eff}}$ at $(p,k)$ is:
      \begin{equation}\label{eq:symbol}
        \sigma_2(L_{\mathrm{eff}})(p,k)
        = -A_\tau k_\tau^2 + A_H
        \Bigl[\Bigl(k_x+\tfrac{y}{2}k_z\Bigr)^2+\Bigl(k_y-\tfrac{x}{2}k_z\Bigr)^2\Bigr].
      \end{equation}
      Following Q5b~\cite{Beau2026q5b} (equations~(17)--(19)), the effective \emph{co-metric
  tensor} $g^{\mu\nu}$ is read off from the principal symbol by $\sigma_2(L_{\mathrm{eff}})(p,k)
  = g^{\mu\nu}(p)k_\mu k_\nu$.
      We denote by $g_{\mu\nu}$ its inverse metric tensor (used for lowering indices and
      computing the temporal projection); these are related by $g^{\mu\rho}g_{\rho\nu}
  = \delta^\mu_\nu$.

      This symbol has rank three: the central slot is empty, and a rank-four extension requires a new operator
      (open problem Q5b-O2), which no source constructs.
      We take a supplied rank-four co-metric of the form of Q11~\cite{Beau2026q11}, Corollary~6.1,
      \begin{equation}\label{eq:metric}
        g^{\mu\nu} = \mathrm{diag}(-A_\tau,\, 2\lambda,\, 2\lambda,\, 2\lambda),
        \qquad A_\tau > 0,\ \lambda > 0,
      \end{equation}
      in an adapted frame.
      It holds under the hypotheses of that corollary: [H-L] and [H-lift], a supplied full-rank extension satisfying
      [T] and [INV$^+$] of Q8~\cite{BeauQ8}, and invariance of its principal symbol under the action [ACT$_\tau$]
      with a negative ordering term.
      Three statements are then distinct.
      First, the invariant form has two independent coefficients: Q8 leaves $\lambda$ free, Q10~\cite{Beau2026q10}
      derives no coefficient value, and Q11 shows that spatial $\mathrm{SU}(2)$ invariance relates $A_\tau$ to
      $\lambda$ by no equation.
      Second, the signature $(-,+,+,+)$ and the absence of mixed temporal--spatial terms are conditional on the
      hypotheses just listed.
      Third, the rescaling $\tau' = \tau\,(2\lambda/A_\tau)^{1/2}$ turns~\eqref{eq:metric} into
      $2\lambda\,\eta^{\mu\nu}$ with $\eta = \mathrm{diag}(-1,+1,+1,+1)$, and the further choice $\lambda = 1$ gives
      the normalisation $g^{\mu\nu} = 2\eta^{\mu\nu}$; this is a choice of units in each sector, not a derived
      value.
      The inverse metric is
      \begin{equation}\label{eq:invmetric}
        g_{\mu\nu} = \mathrm{diag}\bigl(-A_\tau^{-1},\, (2\lambda)^{-1},\, (2\lambda)^{-1},\, (2\lambda)^{-1}\bigr),
      \end{equation}
      so $g_{\tau\tau} = -1/A_\tau < 0$ and $g_{\tau i} = 0$ for the spatial indices~$i$.
      Formulae for the temporal projection in Section~\ref{sec:construction} use
      $g_{\mu\nu}$~\eqref{eq:invmetric}; none of them depends on the values of $A_\tau$ and $\lambda$.

    \subsection{Born--Infeld admissibility budget}

      The quadratic admissibility budget, a structural hypothesis expressing the capacity
      bound [A-cap] of~\cite{Beau2026c}, bounds the total admissible
      capacity increment per projective depth step $dn$:
      \begin{equation}\label{eq:BI}
        (dI_\tau)^2 + (dI_{\mathrm{spat}})^2 = c_{\mathrm{BI}}^2\,dn^2.
      \end{equation}
      After normalisation $c_{\mathrm{BI}} = 1$, the capacity flow vector $\vec{F}_n$
      (Definition~\ref{def:BIsphere} below) lies on the unit Born--Infeld sphere.
      The quadratic equality is a hypothesis of LorCap~\cite{Beau2026n} expressing the inequality [A-cap] of
      the companion paper~\cite{Beau2026c}; neither that paper nor Foundation~\cite{Beau2026m} states it, and it is
      not derived from the principal symbol.

    \subsection{Cone--capacity identification}

      The Born--Infeld budget is normalised by $c_{\mathrm{BI}}$ per depth step, while the co-metric~\eqref{eq:metric}
      has its own light-cone speed $(2\lambda/A_\tau)^{1/2}$ in the units of $\tau$ and of the spatial coordinates.
      Reading a normalised capacity load as the velocity parameter of an observer of the co-metric requires one
      further hypothesis.
      \begin{itemize}
        \item[{[H-cone]}] \emph{Cone--capacity identification.}
        The light-cone speed of the co-metric equals the Born--Infeld bound per ordering step.
        Equivalently, in the rescaled coordinate $\tau'$ in which~\eqref{eq:metric} is proportional to $\eta$, the
        normalised load $\beta = dI_{\mathrm{spat}}/(c_{\mathrm{BI}}\,dn)$ is the coordinate velocity $dx/d\tau'$ of
        the observer it labels, so that $\beta = 1$ lies on the light cone.
      \end{itemize}
      [H-cone] supplies the datum coupling the temporal and spatial normalisations that Q11~\cite{Beau2026q11}
      names as missing (open problem Q11-O1; Q11 lists a saturation rule for the capacity bound among the candidate
      sources).
      It is not derived here or in any source.
      Theorem~\ref{thm:main} does not use it; Corollary~\ref{cor:lorentz} uses it only for its kinematic reading.

  \section{Capacity-Flow Bundle and Born--Infeld Sphere}\label{sec:bundle}

    \subsection{The capacity flow vector}

      In the large-$q$ continuum limit, the cascade depth $n$ is identified with the coordinate
      $\tau$ on $\mathbb{R}_\tau$ through $\tau = n\,\Delta\tau$ with a free step $\Delta\tau$; this is a modelling
      identification (Q11 Modelling identification~3.1~\cite{Beau2026q11}), on the product structure that Q5b
      Remark~3.3~\cite{Beau2026q5b} fixes as a modelling choice.
      The O7 observables become sections of a bundle over $\mathbb{R}_\tau \times \mathrm{Heis}_3(\mathbb{R})$.

      \begin{definition}[Capacity flow vector and unit BI sphere]\label{def:BIsphere}
        The \emph{capacity flow vector} at depth $n$ is
        \[
          \vec{F}_n := F_n^\tau\,\partial_\tau + \vec{F}_n^{\,\mathrm{spat}}
          \;\in\; T_{(\tau,e)}(\mathbb{R}_\tau \times \mathrm{Heis}_3(\mathbb{R})),
        \]
        where $F_n^\tau = dI_\tau/dn$ and $|\vec{F}_n^{\,\mathrm{spat}}| = dI_{\mathrm{spat}}/dn$
        are the normalised temporal and spatial capacity increments.
        The BI budget~\eqref{eq:BI} places $\vec{F}_n$ on the \emph{unit Born--Infeld sphere}:
        \begin{equation}\label{eq:BIsphere}
          (F_n^\tau)^2 + |\vec{F}_n^{\,\mathrm{spat}}|^2 = 1.
        \end{equation}
      \end{definition}

    \subsection{Identification of the spatial component}\label{sec:spatialid}

      The O7 observable $B_n$ is the normalised spatial coverage of the BFS cascade.
      We identify it with the magnitude of the spatial component of the normalised capacity flow:
      \begin{equation}\label{eq:spatial}
        B_n = |\vec{F}_n^{\,\mathrm{spat}}|.
      \end{equation}
      This identification, denoted (S), is a hypothesis of the present paper.
      The Q5b Carnot--Carath\'eodory convergence (Theorem~3.2 of~\cite{Beau2026q5b}) concerns the rescaled BFS
      shells and motivates reading $B_n$ as a spatial footprint; it supplies no capacity-flow vector and does not
      imply~\eqref{eq:spatial}.
      Substituting~\eqref{eq:spatial} into the unit BI sphere~\eqref{eq:BIsphere} gives
      \begin{equation}\label{eq:Ftau}
        F_n^\tau = \sqrt{1 - B_n^2}.
      \end{equation}

  \section{Construction and Uniqueness of $\mathcal{P}_\tau^{\mathrm{Q5b}}$}
    \label{sec:construction}

    \subsection{The linear Lorentzian projector on vectors}

      The completed co-metric~\eqref{eq:metric} endows $\mathbb{R}_\tau \times \mathrm{Heis}_3(\mathbb{R})$
      with a Lorentzian structure.
      The associated \emph{temporal projector} on tangent vectors is the standard Lorentzian
      orthogonal projection onto the timelike direction $\partial_\tau$:
      \begin{equation}\label{eq:Ptau-vector}
        P_\tau V := \frac{g(V,\partial_\tau)}{g(\partial_\tau,\partial_\tau)}\,\partial_\tau
        = \frac{g_{\tau\mu}V^\mu}{g_{\tau\tau}}\,\partial_\tau
        = V^\tau\,\partial_\tau,
      \end{equation}
      using $g_{\tau i} = 0$ from~\eqref{eq:invmetric}, so that $g_{\tau\mu}V^\mu = g_{\tau\tau}V^\tau$.
      The projector does not depend on the values of $A_\tau$ and $\lambda$: it uses only the decoupled temporal
      block and $g_{\tau\tau} \neq 0$.
      It does not use the sign of $g_{\tau\tau}$ either: the Riemannian form
      $\mathrm{diag}(+A_\tau, 2\lambda, 2\lambda, 2\lambda)$ gives the same $P_\tau$.
      The signature makes $\partial_\tau$ timelike; it enters only the kinematic reading of
      Section~\ref{sec:lorentz}.
      The map $V \mapsto P_\tau V = V^\tau\partial_\tau$ is \emph{linear and idempotent}
      on $T(\mathbb{R}_\tau \times \mathrm{Heis}_3(\mathbb{R}))$.

      \begin{remark}[Scope of the linear projector]\label{rem:scope}
        The principal symbol determines $P_\tau$ via the co-metric~\eqref{eq:metric}.
        It does not, however, determine a map on the scalar observable
        $\Sigma_n^{\mathrm{tot}} \in [0,1]$: $P_\tau$ acts on tangent vectors, not on
        real-valued capacity sections.
        The lift from scalars to vectors is provided by the BI bundle structure
        (Definition~\ref{def:BIsphere}), not by the principal symbol.
      \end{remark}

    \subsection{The scalar temporal residual map}

      \begin{definition}[Temporal residual map]\label{def:Ptau}
        The \emph{temporal residual map} $\mathcal{P}_\tau^{\mathrm{Q5b}}$ is the composition:
        \[
          \mathcal{P}_\tau^{\mathrm{Q5b}} \colon
          \Sigma_n^{\mathrm{tot}}
          \xrightarrow{\;\text{BI lift}\;}
          \vec{F}_n
          \xrightarrow{\;P_\tau\;}
          F_n^\tau\,\partial_\tau
          \xrightarrow{\;\text{magnitude}\;}
          |F_n^\tau|.
        \]
        Explicitly, $\mathcal{P}_\tau^{\mathrm{Q5b}}(\Sigma_n^{\mathrm{tot}}) := |F_n^\tau|$,
        where $F_n^\tau$ is the temporal component of the capacity flow vector
        (Definition~\ref{def:BIsphere}).
      \end{definition}

      \begin{remark}[Non-idempotence]\label{rem:nonidempotent}
        The map $\mathcal{P}_\tau^{\mathrm{Q5b}}$, extended by the formula~\eqref{eq:Ptau-scalar} to $[0,1]$, is
        \emph{not} idempotent:
        $(\mathcal{P}_\tau^{\mathrm{Q5b}})^2 \neq \mathcal{P}_\tau^{\mathrm{Q5b}}$ in general.
        It is not an algebraic projection on a space of scalar observables.
        The name ``temporal residual map'' is preferred over ``projector'' for
        the scalar action; the underlying linear map $P_\tau$ on tangent vectors is an
        orthogonal projector.
      \end{remark}

      \begin{remark}[Central direction]\label{rem:central}
        As noted in~\cite{Beau2026n} (Section~5.3), $\mathcal{P}_\tau^{\mathrm{Q5b}}$ does not act
        on the Schr\"odinger representation space $\mathbb{C}^q$.
        The temporal direction is the ordering factor $\mathbb{R}_\tau$, not a direction of
        $\mathrm{Heis}_3(\mathbb{R})$.
        The central direction $\tilde{Z}$ is spatial: its coefficient is the empty central slot of the rank-three
        co-metric (open problem Q5b-O2), and Q8~\cite{BeauQ8} shows that neither the commutator $[X,Y]=Z$ nor any
        lower-order term supplies it.
        The map $\mathcal{P}_\tau^{\mathrm{Q5b}}$ is defined at the level of a supplied completed
        geometry on $\mathbb{R}_\tau\times\mathrm{Heis}_3(\mathbb{R})$, not at the level of
        the Weil representation.
      \end{remark}

    \subsection{Main theorem: explicit formula and uniqueness}

      \begin{theorem}[BI-lifted temporal residual map]\label{thm:main}
        Assume a supplied co-metric of the form~\eqref{eq:metric}, under the hypotheses of Q11 Corollary~6.1
        listed in Section~\ref{sec:cometric}, the identification \textup{(S)} of~\eqref{eq:spatial},
        and the BI capacity-flow lift of Definition~\ref{def:BIsphere}.
        For a map $\psi\colon S_q \to [0,1]$ on the attained values of Section~\ref{sec:objects}, consider, for
        $0 \le n \le R$:
        \begin{itemize}
          \item[\textup{(L)}] $\psi(\Sigma_n^{\mathrm{tot}}) = |F_n^\tau|$, the $\partial_\tau$-coordinate component
          of $P_\tau \vec{F}_n$ (temporal projection of the capacity flow);
          \item[\textup{(BI)}] $B_n^2 + \psi(\Sigma_n^{\mathrm{tot}})^2 = 1$
          (unit BI sphere~\eqref{eq:BIsphere});
          \item[\textup{(W)}] $\Sigma_n^{\mathrm{tot}} = 1 - B_n$ (Weil filling law~\eqref{eq:Sigma}, a proved
          identity that does not involve $\psi$).
        \end{itemize}
        \begin{enumerate}
          \item[\textup{(i)}] \emph{Explicit formula.}
          The temporal residual map $\mathcal{P}_\tau^{\mathrm{Q5b}}$ acts on the attained values
          $\Sigma_n^{\mathrm{tot}} \in S_q$ by
          \begin{equation}\label{eq:Ptau-scalar}
            \mathcal{P}_\tau^{\mathrm{Q5b}}\,\Sigma_n^{\mathrm{tot}}
            \;=\; \sqrt{\Sigma_n^{\mathrm{tot}}\bigl(2-\Sigma_n^{\mathrm{tot}}\bigr)}
            \;=\; \sqrt{1-B_n^2}.
          \end{equation}
          Consequently, the quadratic law
          \begin{equation}\label{eq:quadratic}
            B_n^2 \;+\; \bigl(\mathcal{P}_\tau^{\mathrm{Q5b}}\,\Sigma_n^{\mathrm{tot}}\bigr)^2 \;=\; 1
          \end{equation}
          holds exactly for $0 \le n \le R$.
          \item[\textup{(ii)}] \emph{Scalar uniqueness.}
          $\mathcal{P}_\tau^{\mathrm{Q5b}}$ is the unique map $\psi\colon S_q \to [0,1]$ satisfying \textup{(BI)},
          with $B_n$ expressed by \textup{(W)}.
          Condition \textup{(L)} is not used, and neither is the co-metric.
          A map defined on $[0,1]$ is constrained only through its restriction to $S_q$.
          \item[\textup{(iii)}] \emph{Identification of the component.}
          Under the identification \textup{(S)} and the BI lift, every map satisfying \textup{(L)} satisfies
          \textup{(BI)}, and $\mathcal{P}_\tau^{\mathrm{Q5b}}$ satisfies \textup{(L)}.
          The value fixed in \textup{(ii)} is therefore the $\partial_\tau$-component of the capacity flow, the
          direction $g$-orthogonal to the spatial directions of the supplied co-metric.
        \end{enumerate}
      \end{theorem}

      \begin{proof}
        \textbf{(i) Explicit formula.}
        By the Weil filling law~\eqref{eq:Sigma}: $B_n = 1 - \Sigma_n^{\mathrm{tot}}$.
        By the spatial identification~\eqref{eq:spatial} and the unit BI sphere~\eqref{eq:BIsphere}:
        \[
          (F_n^\tau)^2 = 1 - B_n^2 = 1 - (1-\Sigma_n^{\mathrm{tot}})^2
          = \Sigma_n^{\mathrm{tot}}(2 - \Sigma_n^{\mathrm{tot}}).
        \]
        By Definition~\ref{def:Ptau}: $\mathcal{P}_\tau^{\mathrm{Q5b}}\,\Sigma_n^{\mathrm{tot}}
        = |F_n^\tau| = \sqrt{\Sigma_n^{\mathrm{tot}}(2-\Sigma_n^{\mathrm{tot}})}$.
        Equation~\eqref{eq:quadratic} follows immediately: $B_n^2 + (1-B_n^2) = 1$.
        This identity is exact for $0 \le n \le R$.

        \textbf{(ii) Scalar uniqueness.}
        Since $B_n$ is strictly increasing for $0 \le n \le R$, the map $n \mapsto \Sigma_n^{\mathrm{tot}}$ is
        injective there, so a condition indexed by $n$ is a condition on $\psi$ at the point
        $f = \Sigma_n^{\mathrm{tot}} \in S_q$.
        Let $\psi\colon S_q \to [0,1]$ satisfy (BI).
        By (W): $B_n = 1 - f$.
        Substituting into (BI): $(1-f)^2 + \psi(f)^2 = 1$, hence
        $\psi(f)^2 = 1 - (1-f)^2 = f(2-f)$.
        Since $\psi \ge 0$, $\psi(f) = \sqrt{f(2-f)}$, which is $\mathcal{P}_\tau^{\mathrm{Q5b}}$ on $S_q$ by (i).
        Neither (L) nor any property of the co-metric enters.
        Off $S_q$ none of the three conditions applies, so the values of a map on $[0,1]$ outside $S_q$ are free.

        \textbf{(iii) Identification.}
        If $\psi$ satisfies (L), then $\psi(\Sigma_n^{\mathrm{tot}})^2 = (F_n^\tau)^2 = 1 - B_n^2$ by~\eqref{eq:spatial}
        and~\eqref{eq:BIsphere}, which is (BI).
        $\mathcal{P}_\tau^{\mathrm{Q5b}}$ satisfies (L) by Definition~\ref{def:Ptau}, and $P_\tau$ projects onto
        $\partial_\tau$ along the spatial directions because $g_{\tau i} = 0$ (Section~\ref{sec:construction}).
      \end{proof}

      \begin{remark}[Derivation chain]\label{rem:chain}
        The three conditions play different roles.
        The quadratic law (BI) fixes the scalar values and is necessary for their uniqueness; with the identity (W),
        which converts $B_n$ into a function of the observable, it is sufficient, and (W) gives the values the closed
        form $\sqrt{f(2-f)}$ of the point $f \in S_q$ (Proposition~\ref{prop:necessity}).
        The temporal splitting (L) specifies \emph{which} component of $\vec{F}_n$ is temporal: the
        $\partial_\tau$-component of $P_\tau\vec{F}_n$, with $\partial_\tau$ $g$-orthogonal to the spatial directions.
        It identifies the value fixed by (BI) and (W) as the temporal residual of the capacity flow; it does not
        select that value, and under (S) and the BI lift it implies (BI).
        The Lorentzian kinematic reading, time dilation and the light cone, uses the signature of the co-metric and
        the hypothesis [H-cone] (Corollary~\ref{cor:lorentz}, Remark~\ref{rem:physics}); no step of
        Theorem~\ref{thm:main} uses either.
      \end{remark}

  \section{What Each Input Contributes}\label{sec:necessity}

    \subsection{Which inputs the uniqueness needs}

      \begin{proposition}[Necessity]\label{prop:necessity}
        For the scalar uniqueness of Theorem~\ref{thm:main}\textup{(ii)}:
        \begin{itemize}
          \item the quadratic law is necessary: without it, in the form \textup{(BI)} or in the form of the unit
          sphere~\eqref{eq:BIsphere} entering through \textup{(S)}, the values of $\psi$ on $S_q$ are not fixed;
          \item \textup{(W)} is not needed for the uniqueness of the values $\psi(\Sigma_n^{\mathrm{tot}})$, which
          \textup{(BI)} fixes as $\sqrt{1-B_n^2}$ at each depth; it is needed for their closed form
          $\sqrt{f(2-f)}$ as a function of the point $f \in S_q$;
          \item \textup{(L)} is not necessary: \textup{(BI)} and \textup{(W)} alone fix $\psi$ on $S_q$.
          What \textup{(L)} contributes is the identification of Theorem~\ref{thm:main}\textup{(iii)}.
        \end{itemize}
      \end{proposition}

      \begin{proof}
        Quadratic law: the map $\psi(f) = \sqrt{f}$ on $S_q$ (the candidate $\Sigma_n^B$ of
        Section~\ref{sec:candidates}) takes values in $[0,1]$ and differs from $\sqrt{f(2-f)}$ at every
        $f \in S_q$, since $f \le 1 - 1/q^2 < 1$.
        With $F_n^\tau := \psi(\Sigma_n^{\mathrm{tot}})$ and $|\vec{F}_n^{\,\mathrm{spat}}| := B_n$ it satisfies (L) and
        (S), and $(F_n^\tau)^2 + |\vec{F}_n^{\,\mathrm{spat}}|^2 = 1 - B_n + B_n^2 \neq 1$ for $0 < B_n < 1$, so only the
        quadratic law excludes it.
        Identity (W): at each depth, (BI) alone gives $\psi(\Sigma_n^{\mathrm{tot}}) = \sqrt{1-B_n^2}$, so the values are
        unique without (W).
        The closed form needs (W): if instead $\Sigma_n^{\mathrm{tot}} = g(B_n)$ for another decreasing function $g$,
        (BI) gives $\psi(f) = \sqrt{1 - g^{-1}(f)^2}$, which differs from $\sqrt{f(2-f)}$.
        (W) is a proved identity~\cite{Beau2026n}, so the closed form holds.
        (L): the proof of Theorem~\ref{thm:main}(ii) does not use it; the projector~\eqref{eq:Ptau-vector} uses only
        $g_{\tau i} = 0$ and $g_{\tau\tau} \neq 0$, and the same values follow for either sign of $g_{\tau\tau}$.
      \end{proof}

    \subsection{Comparison with the candidate maps}\label{sec:candidates}

      The note~\cite{Beau2026n} (Section~6) tests three candidate maps
      $\Sigma_n^A$, $\Sigma_n^B$, $\Sigma_n^C$ for $\Sigma_n^\tau$ and finds that all fail the quadratic
      law $B_n^2 + (\Sigma_n^X)^2 \to 1$ with mean deviations of $45\%$, $13\%$, and $7\%$
      respectively.
      Proposition~\ref{prop:failure} states how each differs from the construction.

      \begin{proposition}[Failure modes]\label{prop:failure}
        The three candidate maps of~\cite{Beau2026n} all fail \textup{(BI)}, which is what the numerical test
        of~\cite{Beau2026n} measures.
        They differ from the construction as follows.
        \begin{enumerate}
          \item $\Sigma_n^A = 1 - |{c : \exists(a,b),\,(a,b,c)\in\bigcup_{k\leq n} S_k}|/q$
          (central-coordinate coverage) is not a function of $B_n$ through the BI sphere.
          It measures central BFS coverage, which saturates at depth $n\approx q/2$ while $B_n \approx 1/2$.
          \item $\Sigma_n^B = \sqrt{1-B_n}$ (square-root complement) satisfies the linear budget
          $B_n + (\Sigma_n^B)^2 = 1$ instead of the BI sphere.
          \item $\Sigma_n^C = \sqrt{1-\beta_n^2}$ with $\beta_n = n/n_{\max}$ (depth proxy) has the form of the BI
          sphere, but with the normalised depth $\beta_n$ in place of $B_n$ as spatial component, so it violates the
          identification \textup{(S)}: $\beta_n$ is linear in $n$, while $B_n = (2n^2+2n+1)/q^2$ is quadratic.
        \end{enumerate}
        None of them is excluded by \textup{(W)}, an identity on $\Sigma_n^{\mathrm{tot}}$ that involves no candidate
        for $\Sigma_n^\tau$.
      \end{proposition}

      \begin{proof}
        For $\Sigma_n^A$: the failure is the numerical one of~\cite{Beau2026n} (Table~1); $\Sigma_n^A$ is a function of
        the central coverage, not of $B_n$, so the BI sphere does not constrain it.
        For $\Sigma_n^B$: $B_n^2 + (\Sigma_n^B)^2 = 1 - B_n(1-B_n) \neq 1$ for $0 < B_n < 1$, while
        $B_n + (\Sigma_n^B)^2 = 1$.
        For $\Sigma_n^C$: $B_n^2 + (\Sigma_n^C)^2 = 1$ if and only if $\beta_n = B_n$, and a linear and a quadratic
        function of $n$ agree at no more than two depths.
      \end{proof}

      \begin{remark}[Size of the deviations]\label{rem:info}
        The deviations reported in~\cite{Beau2026n} are numerical.
        For $\Sigma_n^B$ the deviation from the quadratic law is exactly $B_n(1-B_n)$.
        For $\Sigma_n^C$ it is $|\beta_n^2 - B_n^2|$, which depends on the normalisation $n_{\max}$ of the depth proxy;
        \cite{Beau2026n} does not state $n_{\max}$, so no closed-form account of the $7\%$ figure is given here.
      \end{remark}

  \section{The Lorentz Identification}\label{sec:lorentz}

    Theorem~\ref{thm:main} proves Lemma~1 of~\cite{Beau2026n}.
    Theorem~1 of~\cite{Beau2026n} then applies.

    \begin{corollary}[Lorentz factor]\label{cor:lorentz}
      Under the hypotheses of Theorem~\ref{thm:main}, the temporally projected saturation ratio
      \[
        \eta_n^\tau := \frac{B_n}{\mathcal{P}_\tau^{\mathrm{Q5b}}\,\Sigma_n^{\mathrm{tot}}}
        = \frac{B_n}{\sqrt{1-B_n^2}} = \beta\gamma,
      \]
      where $\beta := B_n$ is the normalised spatial mobility and $\gamma = (1-\beta^2)^{-1/2}$.
      Under \textup{[H-cone]}, and with the Lorentzian signature of the co-metric, $\beta$ is the velocity of the
      observer in units of the light-cone speed of the co-metric, and $\gamma$ is its Lorentz factor.
      Consequently,
      \[
        \Phi(\eta_n^\tau) := \bigl(1 + (\eta_n^\tau)^2\bigr)^{-1/2} = \frac{1}{\gamma}.
      \]
    \end{corollary}

    \begin{proof}
      Direct from $\mathcal{P}_\tau^{\mathrm{Q5b}}\,\Sigma_n^{\mathrm{tot}} = \sqrt{1-B_n^2}$
      and the definition of $\gamma$.
    \end{proof}

    \begin{remark}[Physical content]\label{rem:physics}
      The formula of Corollary~\ref{cor:lorentz} holds under the hypotheses of Theorem~\ref{thm:main}.
      Its kinematic reading requires \textup{[H-cone]} and the Lorentzian signature in addition, and then gives:
      (i) \emph{time dilation}: the proper time $\tau_u$ of an observer at load $\beta$, computed with the
      Lorentzian interval of the rescaled co-metric, advances at the rate $d\tau_u/d\tau' = (1-\beta^2)^{1/2}$, which
      equals $F_n^\tau = 1/\gamma$, relative to the rescaled ordering coordinate $\tau'$;
      (ii) \emph{light cone}: on the unit BI sphere a load $\beta = 1$ forces $F^\tau = 0$, i.e.\ no temporal
      residual on the light cone.
      Within the domain of Theorem~\ref{thm:main} the load satisfies $B_n \le B_R = 1/2 + 1/(2q^2)$, so this value
      is not attained by~\eqref{eq:Bn}; it is reached by the exact torus load at depth $n = q-1$, for which the
      quadratic identity holds by the definition of $F^\tau_n$~\cite{Beau2026lco3}.
      Length contraction and the reconstruction of effective Lorentz transformations between
      inertial observers are addressed in~\cite{Beau2026lco1}.
    \end{remark}

  \section{Logical Status}\label{sec:status}

    \subsection{Status table}

      The table below lists the status of every ingredient.

      \smallskip
      \begin{center}
        \renewcommand{\arraystretch}{1.2}
        \begin{tabular}{p{3.4cm}p{5.2cm}p{5.4cm}}
          \hline
          Label & Content & Status \\
          \hline
          \cite{Beau2026n} Cor.~1 & $\Sigma_n^{\mathrm{tot}} = 1-B_n$ exact &
          Proved (Weil factorization) \\
          Q5b Thm.~5.2~\cite{Beau2026q5b} & $L_{\mathrm{eff}}$, rank-three co-metric &
          Conditional on [H-L] (unestablished) and [H-lift] (open) \\
          Q5b-O2 & Rank-four extension & Open; supplied here \\
          Q8~\cite{BeauQ8} & $A_H = A_Z = 2\lambda$, $\lambda$ free &
          Proved for a supplied invariant form ([T], [INV$^+$]); no value of $A_Z$ \\
          Q10~\cite{Beau2026q10} & No coefficient value &
          Ratio bound conditional on [U] and unsupplied inputs \\
          Q11~\cite{Beau2026q11} Cor.~6.1 & Form~\eqref{eq:metric}, $A_\tau$ and $\lambda$ independent &
          Conditional \\
          $g^{\mu\nu} = 2\eta^{\mu\nu}$ & Normalisation after rescaling $\tau$ &
          Choice of units; not used \\
          BI budget~\eqref{eq:BI} & Unit BI sphere &
          Hypothesis of~\cite{Beau2026n} expressing [A-cap] of~\cite{Beau2026c} \\
          (S), \eqref{eq:spatial} & $B_n = |\vec{F}_n^{\,\mathrm{spat}}|$ &
          Identification (this paper) \\
          Definition~\ref{def:BIsphere} & Capacity flow vector $\vec{F}_n$ &
          Structural (this paper) \\
          Theorem~\ref{thm:main}(i) & Formula on $S_q$ &
          Proved under a supplied co-metric~\eqref{eq:metric}, (S) and the BI lift \\
          Theorem~\ref{thm:main}(ii) & Values on $S_q$ fixed by (BI) and (W) &
          Proved; uses neither (L) nor the co-metric \\
          Theorem~\ref{thm:main}(iii) & Value $=$ $\partial_\tau$-component of $P_\tau\vec{F}_n$ &
          Identification under (L), (S) and the BI lift; signature not used \\
          Proposition~\ref{prop:necessity} & Quadratic law necessary for the uniqueness, (W) for the closed form;
          (L) not &
          Proved \\
          Corollary~\ref{cor:lorentz} & $\Phi(\eta_n^\tau) = 1/\gamma$ &
          Formula proved under the same hypotheses; kinematic reading conditional on [H-cone] and the signature \\
          {[H-cone]} & Cone speed $=$ BI bound per ordering step &
          Hypothesis; the coupling datum of Q11-O1, open \\
          \hline
        \end{tabular}
      \end{center}

      \medskip
      The analytic content is conditional on a supplied rank-four co-metric of the form~\eqref{eq:metric}, which in
      turn rests on the hypotheses of Q11 Corollary~6.1, and on the identification (S).
      No coefficient value enters: the projector $P_\tau$ uses only the decoupled temporal block and
      $g_{\tau\tau} \neq 0$; the signature enters only the kinematic reading.
      The spectral universality hypothesis [U] enters no step of the proof; it is open, U1~\cite{Beau2026u1}
      recording that no proof of it is available.
      The BI lift (Definition~\ref{def:BIsphere}) uses the quadratic budget hypothesis~\eqref{eq:BI} and the
      identification (S).
      The kinematic reading of Corollary~\ref{cor:lorentz} uses [H-cone] and the signature in addition.

    \subsection{Status of LC-O1--O3}

      \paragraph{LC-O1 (Length contraction).}
        Effective Lorentz transformations between inertial observers are reconstructed
        in~\cite{Beau2026lco1}.
        The Born--Infeld capacity sphere supplies the load $\beta$ and the residual
        $F^\tau = 1/\gamma$, a supplied co-metric of the form~\eqref{eq:metric}, rescaled in $\tau$, supplies the
        Lorentzian interval whose isometries are the Lorentz group, and [H-cone] identifies the load with the
        velocity of the boost.
        Length contraction follows as the reciprocal spatial projection onto the
        simultaneity hyperplane of the moving splitting.

      \paragraph{LC-O2-O1 (Capacity--metric bridge).}
        LCII~\cite{Beau2026lco2} postulates the equivalence of the capacity definition
        $R(x)^2 = 1-\epsilon(x)$ and the metric definition $R(x)^2 = g_{\tau\tau}(x)/g_{\tau\tau}^\infty$, whose
        flat-space reference value is postulated.
        LC-O2-O1~\cite{Beau2026lco2o1} proves that the scalar free-fraction reduction does
        not establish this equivalence, so the bridge remains open.

      \paragraph{LC-O3 (Finite-$q$ corrections at saturation).}
        Finite-$q$ corrections near saturation are analysed in~\cite{Beau2026lco3}.
        The quadratic identity $B_n^2 + (F^\tau_n)^2 = 1$ is exact for every $n$ with
        the exact torus load.
        The corrections affect the parabolic Carnot identification
        $B_n^{\mathrm{parab}} = (2n^2+2n+1)/q^2$, not the capacity identity itself:
        for $n = R+m$,
        \[
          B_n^{\mathrm{torus}} = B_n^{\mathrm{parab}} - \frac{4m^2}{q^2}.
        \]
        The parabolic identification is exact for $n \leq R := (q-1)/2$ and
        asymptotically exact for $n = R + o(q)$.

  \section{Conclusion}
  \label{sec:conclusion}

    The quadratic Born--Infeld law and the Weil filling law fix the scalar temporal residual at the attained values
    of the total residual, $F_n^\tau=\sqrt{1-B_n^2}$.
    The temporal splitting of the supplied co-metric identifies that value as the temporal component of the
    capacity flow; it does not select it.
    Under [H-cone], with the Lorentzian signature, the value is the standard time-dilation factor $1/\gamma$.
    The splitting uses a supplied co-metric of the conditional form of Q11 with two independent coefficients; no
    coefficient value is used, and the normalisation $g^{\mu\nu} = 2\eta^{\mu\nu}$ is a choice of units.
    It does not close the non-homogeneous capacity--metric bridge, for which the scalar
    free-fraction construction has a proved no-go.

    \paragraph{Interpretive outlook.}
    Under [H-cone], the result supports viewing relativistic clock-rate reduction as residual capacity after
    spatial mobility has been accounted for.
    This is a structural reading of the proved map, not an independent dynamical law.

  \section*{Acknowledgements}

  The author acknowledges the use of large language models as a supportive tool for refining
  language, structure, and internal consistency during the development of this manuscript.

  \newpage

  \bibliographystyle{plain}
  \bibliography{cosmochrony-bibliography}

\end{document}
